\section{高斯解码细节}
\label{sec:supp-codec}

每个高斯存储中心 $\boldsymbol\mu$、协方差 $\boldsymbol\Sigma$、不透明度 $\alpha$ 和方向外观。

背景、刚体和人体高斯解码器分别将每个查询独立映射为 $K=4$ 个子高斯。每个子高斯输出 23 个参数，包括三个中心偏移、三个尺度、四个旋转参数、一个不透明度和十二个外观系数。查询 $i$ 的子高斯 $k$ 的中心被限制在其支撑单元内：

\begin{equation}
\boldsymbol\mu_{i,k}=\mathbf q_i+
\tfrac12\boldsymbol\delta_i\odot\tanh(\mathbf b_{i,k}).
\label{eq:child-center}
\end{equation}

正尺度与归一化旋转确定正定协方差。不透明度乘以 $1-f_i$，使自由空间查询不能产生可见高斯。人体子高斯继承支撑查询的蒙皮权重；刚体子高斯共用所属资产的刚体变换。独立天空头预测立方体贴图颜色，并在 $V_{\rm sky}$ 中保留观测支持的方向域。

外观为每个颜色通道存储常数 RGB 项和三个局部方向系数，按 $c(d)=\max(0,c_0+C_{\rm app}\hat d_{\rm local})$ 求值。解码逐查询执行：编解码器冻结后，不变的已见潜码产生不变的高斯参数。补全网络可以在查询之间交换上下文，同时保留这一解码性质。
